% Generated draft: requires mathematical and source review.
\documentclass[11pt]{article}
\usepackage[margin=25mm]{geometry}
\usepackage{fontspec}
\setmainfont{Latin Modern Roman}
\usepackage{amsmath,amssymb,mathtools,unicode-math}
\setmathfont{Latin Modern Math}
\usepackage{enumitem,needspace,xurl,hyperref}
\setlength{\parindent}{0pt}
\setlength{\parskip}{5pt}
\setlength{\emergencystretch}{3em}
\newcommand{\N}{\mathbb{N}}
\newcommand{\Z}{\mathbb{Z}}
\newcommand{\Q}{\mathbb{Q}}
\newcommand{\R}{\mathbb{R}}
\newcommand{\C}{\mathbb{C}}
\newcommand{\F}{\mathbb{F}}
\newcommand{\A}{\mathbb{A}}
\newcommand{\PP}{\mathbb{P}}
\newcommand{\E}{\mathbb{E}}
\newcommand{\eps}{\varepsilon}
\newcommand{\dd}{\,\mathrm{d}}
\begin{document}
\section*{UnsolvedMath ID: 701143}
\textbf{Source problem code:} catalogue entry
\subsection{Definitions and mathematical statement}
Let \(M\), \(B\), and \(F_{24}'\) denote the Monster, Baby Monster, and sporadic simple group denoted by \(F_{24}'\) in the supplied record, respectively.

For a finite group \(G\), write \(o(g)\) for the order of \(g\in G\). A conjugacy class \(T\) of involutions makes \(G\) a \(6\)-transposition group if
\[
G=\langle T\rangle,\qquad o(st)\leq 6\quad\text{for every }s,t\in T.
\]
For involutions \(s,t\), the subgroup \(\langle s,t\rangle\) is dihedral of order \(2o(st)\), allowing the degenerate case \(s=t\), when this subgroup has order \(2\). Define
\[
\mathcal D(G,T)=\{\langle s,t\rangle:s,t\in T\}/G,
\]
where the quotient denotes conjugacy of subgroups under \(G\). For \([D]\in\mathcal D(G,T)\), put \(n([D])=|D|/2\); this equals \(o(st)\) for any generating pair of involutions defining \(D\).

To specify a possible meaning of the root coefficients, let \(\Phi\) be a finite irreducible root system of type \(E_r\), choose simple roots \(\alpha_1,\ldots,\alpha_r\), and write its highest root as
\[
\theta=\sum_{i=1}^{r}a_i\alpha_i.
\]
The marks of the associated untwisted affine diagram are \(a_0=1,a_1,\ldots,a_r\). Equivalently, with affine simple root \(\alpha_0=\delta-\theta\), they satisfy
\[
\delta=\sum_{i=0}^{r}a_i\alpha_i.
\]
This gives a conventional interpretation of the source's phrase ``coefficient of the highest root of affine type''; an affine root system does not itself have a highest root in the finite-root-system sense.

\textbf{(a) Explanatory problem.} The source asserts that \(M\) is a \(6\)-transposition group and that pairs of its ``Fischer transpositions'' generate nine conjugacy classes of dihedral subgroups. Let \(T_M\) provisionally denote the intended class of involutions. Under the affine-mark interpretation above, the proposed numerical correspondence is
\[
|\mathcal D(M,T_M)|=9,
\qquad
\bigl\{\!\bigl\{n([D]):[D]\in\mathcal D(M,T_M)\bigr\}\!\bigr\}
=
\bigl\{\!\bigl\{a_0,a_1,\ldots,a_8\bigr\}\!\bigr\}
\quad(\Phi=E_8),
\]
where double braces denote a multiset. Explain the relationship between these group-theoretic data and the affine \(E_8\) diagram.

The source additionally asks for analogous explanations relating \(B\) to \(E_7\) and \(F_{24}'\) to \(E_6\), using the central covers denoted by \(\widetilde B=2.B\) and \(\widetilde F=3.F_{24}'\), with products read modulo their centres. Explicitly, for a specified pair \(x,y\in\widetilde G\), the relevant product order under this instruction is
\[
o\bigl((xZ(\widetilde G))(yZ(\widetilde G))\bigr)
=o\bigl(xyZ(\widetilde G)\bigr)
\quad\text{in }\widetilde G/Z(\widetilde G).
\]
The source does not specify the corresponding element classes, admissible pairs, or exact correspondence with the \(E_7\) and \(E_6\) diagrams. Consequently these analogies remain explanatory targets requiring clarification, rather than fully quantified assertions supplied by the record. The class \(T_M\), the treatment of equal pairs, and the multiset interpretation above also require confirmation against the original source.

\textbf{(b) Explanatory problem.} For a finite group \(G\), let \(\operatorname{Mult}(G)=H_2(G,\Z)\) be its Schur multiplier. For a finite root system \(\Phi\), let \(Q(\Phi)\) be the lattice generated by its roots and
\[
P(\Phi)=\{\lambda\in\operatorname{span}_{\R}(\Phi):
\langle\lambda,\alpha^\vee\rangle\in\Z\text{ for every }\alpha\in\Phi\}
\]
its weight lattice, where \(\alpha^\vee=2\alpha/(\alpha,\alpha)\). A conventional meaning of the fundamental group of a root system is \(P(\Phi)/Q(\Phi)\). Under this interpretation, the observation to be explained is the sequence of abstract group isomorphisms
\[
\begin{aligned}
\operatorname{Mult}(M)&\cong P(E_8)/Q(E_8)\cong\{0\},\\
\operatorname{Mult}(B)&\cong P(E_7)/Q(E_7)\cong\Z/2\Z,\\
\operatorname{Mult}(F_{24}')&\cong P(E_6)/Q(E_6)\cong\Z/3\Z.
\end{aligned}
\]
Explain why these correspondences occur. The source does not identify a particular Lie group or specify the meaning of ``fundamental group''; the lattice interpretation is therefore explicit and provisional. Neither part specifies a formal criterion for what would constitute an explanation. These are source-reported observations and explanatory questions, not mathematical conclusions established by this draft.
\subsection{Short English statement}
Explain the reported relationships between the Monster, Baby Monster, and \(F_{24}'\) and the root systems \(E_8,E_7,E_6\): first through orders of products of distinguished elements and affine root coefficients, and second through Schur multipliers and root-system fundamental groups. The precise element classes, the two smaller-group analogies, and the intended meaning of an explanation require clarification. The registry labels the problem open; this draft does not independently verify that status.
\subsection{Sources}
\begin{itemize}
\item Local UnsolvedMath registry: permanent ID \(701143\), exact source problem code \texttt{catalogue entry}, title ``Kourovka Notebook Problem 15.55.'' The supplied record attributes the problem to J. McKay, proposed in 2002, and labels it \texttt{open}.
\item \emph{The Kourovka Notebook (Unsolved Problems in Group Theory)}, original problem 15.55, 15th issue (2002). The supplied background also identifies the 21st-issue edition (2026), arXiv:1401.0300v48: \url{https://arxiv.org/abs/1401.0300v48}. The supplied collection description names E. I. Khukhro and V. D. Mazurov as editors. These edition details are supplied attribution and have not been independently verified.
\item J. McKay, contribution in \emph{Finite groups, Santa Cruz Conf. 1979}, \emph{Proc. Symp. Pure Math.}, volume 37, American Mathematical Society, Providence, RI, 1980, pp. 183--186. The supplied record gives no contribution title.
\item R. E. Borcherds, \emph{Doc. Math., J. DMV Extra Vol. ICM Berlin 1998}, volume I (1998), pp. 607--616. The supplied record gives no article title.
\item C. H. Lam, H. Yamada, and H. Yamauchi, \emph{Trans. Amer. Math. Soc.} \textbf{359}, no. 9 (2007), pp. 4107--4123. The source includes an ``Editors' Comment of 2005'' reporting that some progress was made in this work. Both the comment's stated year and the publication year are retained as supplied; the scope of that progress has not been independently verified. No article title is supplied.
\end{itemize}
\end{document}
